\section{Conclusões}

\begin{frame}{Conclusões e trabalhos futuros}
  \begin{itemize}[<+->]
    \setlength{\itemsep}{0.1cm}
    \item EADRC para supressão do tremor de punho na DP, considerando a dinâmica acoplada do membro superior (modelo de 3 GdL).
    \item Alta supressão com menor esforço de controle que o DE-PID; entre as variantes EADRC, o ZPLP gerou sinais de controle mais suaves.
    \item EADRC+ZPLP: TPSR médio de 61,31\% em movimento, superior aos controladores clássicos realizáveis.
    \item Desempenho consistente com variações de até 50\% em rigidez e amortecimento: adequado quando as propriedades biomecânicas são incertas ou variantes no tempo.
    \item A estimação do movimento voluntário é decisiva para não deixar tremores de menor frequência fora da perturbação total estimada.
  \end{itemize}

  \pause
  \vspace{0.2cm}
  \begin{exampleblock}{Trabalhos futuros}
    Outras estratégias de controle da literatura, métodos alternativos de estimação do movimento voluntário e outros modelos do membro superior.
  \end{exampleblock}
\end{frame}
